\section{Cosmology and the Predictions}
\label{sec:cosmology}

A universe that grows forever has a cosmology built in. This section walks the cosmological findings
and the sharp observational predictions (P13, P30, P46, P26, P27, P78, P82).

\paragraph{Expansion and inflation (P13, P30).}
P13 gives the cosmology: a growing mesh naturally expands, with an arrow of time from the one-way
growth. P30 shows an inflationary phase, a brief early burst of fast growth that then settles, from
the growth rule itself.

\paragraph{The everpresent cosmological constant (P46).}
P46 gives the model's most striking quantitative contact. The cosmological constant is not a tuned
number but the natural fluctuation of the spacetime volume, scaling as one over the square root of
that volume. Evaluated at the size of our universe, it lands near the observed, famously tiny,
dark-energy scale, to order of magnitude, with no fine-tuning. This is the everpresent or dynamical
$\Lambda$ of the causal-set programme, realised on the substrate.

\paragraph{The dark sector.}
Dark energy is the everpresent $\Lambda$ above. Dark matter has a candidate mechanism, a small
long-range correction to gravity that flattens galaxy rotation curves without a new particle, part
of the field-content work of P18 to P22.

\paragraph{The primordial seed (P78).}
P78 is the first step on structure formation. A sprinkled substrate is a Poisson process, so the
density contrast in a region falls as one over the square root of the count it holds, the same
one-over-root-volume law as the everpresent constant. The measured exponent is $-0.509$, against the
Poisson prediction $-1/2$, a clean scale-free seed. Turning that seed into the observed, slightly
tilted spectrum of the microwave background, through inflationary stretching and gravitational
growth, is the larger task that remains.

\paragraph{Sharp predictions against the bounds (P26, P27, P82).}
The model's two distinctive signatures are the swerve and Lorentz safety. The swerve (P26) is a tiny
momentum diffusion from discreteness: a particle cannot keep a perfectly straight trajectory on a
discrete substrate, so its rapidity random-walks, with a variance growing linearly in proper time.
The diffusion rate falls as the discreteness fines (measured scaling $\sim$ density$^{-1.6}$), so at
the Planck scale it is far below the cosmic-ray bound. Lorentz safety (P27) is the geometric result
of Section~\ref{sec:geometry}: the substrate predicts no first-order Lorentz violation. P82 sets
these against the tightest current bounds. The model passes the Fermi-LAT gamma-ray-burst timing
bound on linear Lorentz violation (coefficient below $1/7.6$), and the test is discriminating, since
the same bound excludes a regular lattice, whose photon speed becomes order-one anisotropic near the
cutoff. The sharp, falsifiable claim is the Lorentz one: a confirmed first-order energy-dependent
photon speed would falsify the Lorentz-safe substrate.

